\itemhead{1.2 思考 · 未编号题组（教材p.20）}{XJ-B1-C01-TH-P020-Y0942}
\label{XJ-B1-C01-TH-P020-Y0942}
{\footnotesize 湘教版教材p.20 / PDF第27页 / 1.2 思考 / 未编号题组}
\field{原题}{如何判断全称量词命题的真假？如何判断存在量词命题的真假？}
\field{参考答案}{全称命题真须对范围内每一个对象证明，假只需一个反例；存在命题真只需给出一个对象，假须证明范围内所有对象均不满足。}
{\small \field{关联知识}{L03 全称、存在量词及真假的判定（教材帮p.38、39 / PDF47、48）}
\field{方法与易错点}{先确认量词、作用对象及论域，再采用任意性证明或实例反例。 易错：用几个满足的数证明全称命题；把“存在”理解为“唯一存在”。}
\field{答案依据}{答案册未提供本栏目的独立对应答案；依据原题、教材概念补写参考回应。开放性总结不设唯一数值答案。}}
